\documentclass{ujarticle}
\usepackage[utf8]{inputenc}
\usepackage[T1]{fontenc}
\usepackage{geometry}
\usepackage{graphicx}
\usepackage{float}
\usepackage{hyperref}
\usepackage[english,indonesian]{babel}

\title{Optimasi Pemrosesan Template LaTeX Menggunakan Pendekatan Berbasis AI}
\author{Academic Support Agent}
\date{\today}

\begin{document}
\maketitle

\begin{abstract}

\end{abstract}

\chapter{PENDAHULUAN}


\chapter{TINJAUAN PUSTAKA}


% Original content (commented out):
% %\documentclass[11pt,a4paper,uplatex,dvipdfmx]{ujarticle} 		% for uplatex
\documentclass[11pt,a4j,dvipdfmx]{jarticle} 					% for platex
\input{pieces/form00_header} % pieces
\input{pieces/kakenhi7} % pieces
\input{pieces/form01_header} % pieces
\input{pieces/form02_2024_header} % pieces
\input{pieces/hook3} % pieces
%#Name: ilr-1
\input{pieces/form04_jsps_eng_headers} % pieces
\input{pieces/form04_ilr-1_header} % pieces
% ===== Global definitions for the Kakenhi form ======================
%　基本情報
%
%------ 研究課題名  -------------------------------------------
\newcommand{\研究課題名}{象の卵}

%----- 研究機関名と研究代表者の氏名-----------------------
\newcommand{\研究機関名}{逢坂大学}
\newcommand{\研究代表者氏名}{湯川秀樹}
\newcommand{\me}{\underline{\underline{H.~Yukawa}}} 
\input{pieces/inst_general_images} % pieces
% user07_header
% ===== my favorite packages ====================================
% ここに、自分の使いたいパッケージを宣言して下さい。
\usepackage{wrapfig}
%\usepackage{amssymb}
%\usepackage{mb}
%\DeclareGraphicsRule{.tif}{png}{.png}{`convert #1 `dirname #1`/`basename #1 .tif`.png}
\usepackage{lineno}

% ===== my personal definitions ==================================
% ここに、自分のよく使う記号などを定義して下さい。
\newcommand{\klpionn}{K_L \to \pi^0 \nu \overline{\nu}}
\newcommand{\kppipnn}{K^+ \to \pi^+ \nu \overline{\nu}}

% ----- 業績リスト用 -------------
\newcommand{\paper}[6]{%
	% paper{title}{authors}{journal}{vol}{pages}{year}
	\item ``#1'', #2, #3 {\bf #4}, #5 (#6).			% お好みに合わせて変えてください。
}

\newcommand{\etal}{\textit{et al.\ }}
\newcommand{\ca}[1]{*#1}	% corresponding author;   \ca{\yukawa}  みたいにして使う
\newcommand{\invitedtalk}{招待講演}

\newcommand{\yukawa}{H.~Yukawa}					% no underline
%\newcommand{\yukawa}{\underline{\underline{H.~Yukawa}}}	% with 2 underlines
\newcommand{\tomonaga}{S.~Tomonaga}

\newcommand{\prl}{Phys.\ Rev.\ Lett.\ }		% よく使う雑誌も定義すると楽

% ===== 欄外メモ ==================
\newcommand{\memo}[1]{\marginpar{#1}}
%\renewcommand{\memo}[1]{}	% 全てのメモを表示させないようにするには、行頭の"%"を消す

\input{pieces/hook5} % pieces

\begin{document}
\input{pieces/hook7} % pieces
%#Split: 01_collaboration_abstract  
%#PieceName: p01_collaboration_abstract
\input{pieces/p01_collaboration_abstract_00}
\section{国際共同研究の枠組みと意義の概要}
%　　　　＜＜最大　1ページ＞＞

%s01_collab_abstract_ilr
\JSPSInstructions		% <-- 留意事項。これは消すか、コメントアウトしてください。
%begin 国際的研究の意義と必要性の概要 ＝＝＝＝＝＝＝＝＝＝＝＝＝＝＝＝＝＝＝＝
	象の生息地に行かないと真の研究はできない。
	象牙の価値を下げないと、象の生息地から象が消える。
	象牙の価値を下げる最良の方法は、国際協力で象の卵の殻から象牙の代替品を作ることである。
%end 国際的研究の意義と必要性の概要 ＝＝＝＝＝＝＝＝＝＝＝＝＝＝＝＝＝＝＝＝

\input{pieces/p01_collaboration_abstract_01}

%#Split: 02_research_abstract  
%#PieceName: p02_research_abstract
\input{pieces/p02_research_abstract_00}
\section{研究計画の概要}
%　　　　＜＜最大　1ページ＞＞

%s01_research_abstract_ilr
%begin 研究計画の概要 ＝＝＝＝＝＝＝＝＝＝＝＝＝＝＝＝＝＝＝＝
	象の卵の研究計画と方法は．．．
	\vspace*{10zw}	% （概要）と（本文）の間が10行程度になるよう、必要に応じて値を調整してください。	
%end 研究計画の概要 ＝＝＝＝＝＝＝＝＝＝＝＝＝＝＝＝＝＝＝＝

\input{pieces/p02_research_abstract_01}

%#Split: 03_education_abstract  
%#PieceName: p03_education_abstract
\input{pieces/p03_education_abstract_00}
\section{人材育成計画の概要}
%　　　　＜＜最大　1ページ＞＞

%s01_edu_abstract_ilr
%begin 人材育成計画の概要 ＝＝＝＝＝＝＝＝＝＝＝＝＝＝＝＝＝＝＝＝
	時間がかかり、非効率的に思われるかもしれないが、
	もっとも効果的な人材育成の方法は、若い人に研究を任せきることである。
	人は失敗してこそ、真に学ぶ。
	失敗は計画できない。
	したがって、人材育成の計画はできない。
%end 人材育成計画の概要 ＝＝＝＝＝＝＝＝＝＝＝＝＝＝＝＝＝＝＝＝
\input{pieces/p03_education_abstract_01}

%#Split: 99_tail
\input{pieces/hook9} % pieces
\end{document}



\end{document}